\documentclass[10pt,letterpaper]{article}
\usepackage[utf8]{inputenc}
\usepackage[T1]{fontenc}
\usepackage{amsmath,amssymb,amsfonts}
\usepackage{booktabs}
\usepackage{multirow}
\usepackage{graphicx}
\usepackage{xcolor}
\usepackage[margin=1in]{geometry}
\usepackage{hyperref}
\usepackage{natbib}
\usepackage{caption}
\usepackage{subcaption}
\usepackage{array}

\newcommand{\ie}{\textit{i.e.}}
\newcommand{\eg}{\textit{e.g.}}
\newcommand{\cf}{\textit{cf.}}
\newcommand{\etc}{\textit{etc.}}

\title{Diagnosing Conditioning-Induced Value Overestimation\\in Variable-Rate World-Model Imagination}

\author{
  Anonymous Authors
}

\date{}

\begin{document}
\maketitle

%==============================================================================
\begin{abstract}
%==============================================================================
Temporal abstraction in world-model reinforcement learning promises improved planning efficiency through variable-rate imagination.
We present a systematic diagnostic analysis of variable-rate imagination in DreamerV3, revealing that trained temporal embeddings (\texttt{dt\_embedding}) reliably trigger \emph{conditioning-induced value overestimation} (CIVO)---a pathological feedback loop where the critic exploits precise but inaccurate temporal conditioning to produce systematically inflated value estimates.
Across two DeepMind Control Suite environments (Walker Walk, Cheetah Run) and multiple random seeds, we find that: (1) all variable-rate imagination configurations degrade performance relative to dt=1 baselines; (2) trained \texttt{dt\_embedding} reliably induces normalized value-to-eval (NVE) ratios exceeding $1.5\times$ baseline, serving as a diagnostic marker with 83\% classification accuracy across environments; (3) removing or freezing the temporal embedding eliminates CIVO but introduces distinct failure modes (reward degradation or gradient starvation); and (4) restricting imagination to dt=1 eliminates the CIVO feedback loop (Pathway~A) but does not reduce overall collapse probability---a distinct policy instability pathway (Pathway~B) manifests at the same 50\% rate across 4 seeds.
A sensitivity analysis estimates CIVO's independent contribution to performance degradation at 20--35 percentage points, depending on decomposition assumptions.
These findings reveal that the vulnerability is architectural: the trained temporal embedding itself, not the imagination strategy, is the root cause.
Our results provide actionable guidance for future temporal abstraction designs in imagination-based actor-critic architectures.
\end{abstract}

%==============================================================================
\section{Introduction}
\label{sec:intro}
%==============================================================================

Temporal abstraction---the ability to reason and act at multiple time scales---is a longstanding goal in reinforcement learning~\citep{sutton1999options,vezhnevets2017feudal}.
World-model RL offers a natural testbed for temporal abstraction: by modifying the temporal stride (\texttt{dt}) of imagination rollouts within a learned dynamics model, agents can plan over variable time horizons without altering environment interaction.
DreamerV3~\citep{hafner2023dreamerv3} provides a well-characterized platform for such investigations, with its Recurrent State-Space Model (RSSM) serving as the dynamics backbone.

Despite the appeal, our systematic exploration of variable-rate imagination in DreamerV3 reveals a consistent negative result: \emph{all} configurations with dt$>$1 degrade performance relative to the dt=1 baseline.
This finding contrasts with recent successes of temporal conditioning in model-predictive control settings~\citep{nhu2025tawm}, suggesting an architecture-specific vulnerability in imagination-based actor-critic training.

We identify the root cause as \textbf{conditioning-induced value overestimation (CIVO)}: a feedback loop in which a trained temporal embedding enables the critic to produce precise but systematically inflated value estimates at dt$>$1, which the actor then exploits, leading to catastrophic performance collapse.
Crucially, we find that this vulnerability operates through \emph{multiple pathways}: eliminating the value overestimation loop (Pathway~A) via dt=1 imagination does not prevent collapse---a distinct policy instability pathway (Pathway~B) emerges at the same rate, driven by the trained temporal embedding's presence in the training data.

\paragraph{Contributions.}
\begin{itemize}
    \item A factorial experiment design isolating trained \texttt{dt\_embedding} as the necessary condition for CIVO, validated across two DMC environments and multiple seeds.
    \item The normalized value-to-eval (NVE) ratio as a diagnostic marker: NVE~$>1.5\times$ achieves 83\% classification accuracy for identifying CIVO-vulnerable configurations (100\% sensitivity, 83\% specificity).
    \item Discovery of dual collapse pathways: Pathway~A (classic CIVO, value-driven, entropy decreasing) and Pathway~B (policy instability, eval-driven, entropy increasing), both triggered by trained \texttt{dt\_embedding}.
    \item A complete ablation matrix (trained/untrained/frozen/absent \texttt{dt\_embedding}) characterizing four distinct failure modes.
    \item Quantification of CIVO's independent contribution at 20--35 percentage points via sensitivity analysis.
\end{itemize}

%==============================================================================
\section{Background and Related Work}
\label{sec:related}
%==============================================================================

\paragraph{Temporal Abstraction in RL.}
The options framework~\citep{sutton1999options} formalized temporal abstraction through temporally extended actions.
Subsequent work explored learned temporal abstractions including FiGAR~\citep{sharma2017figar}, action repeat~\citep{lakshminarayanan2017dynamic}, and hierarchical policies~\citep{vezhnevets2017feudal,nachum2018hiro}.
These approaches modify the \emph{action} space; in contrast, we study temporal abstraction in the \emph{imagination} space of world models.

\paragraph{World-Model RL.}
The Dreamer family~\citep{hafner2020dreamerv1,hafner2021dreamerv2,hafner2023dreamerv3} learns latent dynamics models (RSSM) and trains actor-critic policies entirely within imagined rollouts.
Temporal abstraction in this setting means varying the time stride of imagination steps.
TAWM~\citep{nhu2025tawm} successfully uses dt-conditioning in TD-MPC2's model-predictive control setting, where value estimates are not propagated through long imagination chains.
Our finding that the same mechanism fails in imagination-based actor-critic training highlights an architecture-specific vulnerability.

\paragraph{Value Overestimation in MBRL.}
Model exploitation and value overestimation are known failure modes in model-based RL~\citep{janner2019mbpo,chua2018pets}.
Lu et al.~\citep{lu2018nondelusional} formalize ``delusional bias'' arising from policy-class constraints.
Our CIVO mechanism is distinct: it arises from the interaction between learned temporal conditioning and model prediction errors, not from policy constraints.

%==============================================================================
\section{Method}
\label{sec:method}
%==============================================================================

\subsection{Base Architecture}

We build on DreamerV3~\citep{hafner2023dreamerv3} with its RSSM world model, actor-critic imagination training, and symlog predictions.
The RSSM learns latent dynamics $s_{t+1} = f_\theta(s_t, a_t)$ and observation/reward decoders.
The actor and critic are trained entirely within $H$-step imagination rollouts from replay buffer states.

\subsection{Variable-Rate Imagination}

We introduce temporal abstraction through three mechanisms:
\begin{enumerate}
    \item \textbf{Temporal embedding} (\texttt{dt\_embedding}): A learned embedding $e(\Delta t)$ concatenated to the RSSM state, conditioning all predictions on the temporal stride.
    \item \textbf{Temporal subsampling}: When enabled, training transitions are subsampled at stride $\Delta t$, so each imagination step covers $\Delta t$ environment steps.
    \item \textbf{Imagination dt}: The $\Delta t$ value used during actor-critic imagination rollouts (may differ from training data).
\end{enumerate}

\subsection{Factorial Experiment Design}

We design a controlled factorial grid varying:
\begin{itemize}
    \item \textbf{dt\_embedding state}: trained / untrained (random fixed) / frozen (random, \texttt{requires\_grad=False}) / absent (\texttt{dim=0})
    \item \textbf{Temporal subsampling}: ON (dt=2 transitions in training) / OFF (dt=1 only)
    \item \textbf{Imagination horizon}: $H=8$ / $H=15$
    \item \textbf{Imagination dt}: 1 or 2
\end{itemize}

\subsection{Diagnostic Metrics}

\paragraph{Normalized Value-to-Eval Ratio (NVE).}
We define the NVE as the ratio of experiment value-to-eval (v/e) to baseline v/e at the same training step:
\begin{equation}
    \text{NVE}(t) = \frac{v_{\text{exp}}(t) / e_{\text{exp}}(t)}{v_{\text{base}}(t) / e_{\text{base}}(t)}
    \label{eq:nve}
\end{equation}
NVE $>1$ indicates the experiment's critic is more optimistic relative to its own performance than the baseline's critic is to its performance.
We find that sustained NVE $>1.5\times$ at steps $\geq$35K reliably indicates pathological overestimation.

\subsection{Environments}
Experiments are conducted on DeepMind Control Suite~\citep{tassa2018dmcontrol} Walker Walk and Cheetah Run.
Cartpole Swingup was planned but not completed within the compute budget.

%==============================================================================
\section{Experiments}
\label{sec:experiments}
%==============================================================================

\subsection{All Variable-Rate Configurations Degrade Performance (C1)}

Table~\ref{tab:c1_degradation} summarizes the key result: every dt$>$1 imagination configuration degrades relative to the dt=1 baseline ($\sim$950 eval return in Walker Walk).

\begin{table}[h]
\centering
\caption{Variable-rate imagination consistently degrades performance in Walker Walk. Baseline eval $\approx$916--950.}
\label{tab:c1_degradation}
\begin{tabular}{lcc}
\toprule
Configuration & Peak Eval & Outcome \\
\midrule
Baseline (dt=1) & $\sim$950 & Stable convergence \\
Fixed dt=2, subsample ON, H=15 & 485 @35K & Collapse to $\sim$135 @65K \\
Fixed dt=4 & 177 @25K & Severe degradation (33\% of baseline) \\
Learned gate & 541 @55K & Collapse to 99 @145K \\
\bottomrule
\end{tabular}
\end{table}

Increasing dt from 2 to 4 worsens degradation monotonically (peak eval: dt=1 $\sim$950 $>$ dt=2 485 $>$ dt=4 177), consistent with RSSM prediction error compounding over longer temporal gaps.

\subsection{Trained dt\_embedding Triggers Pathological Overestimation (C2)}
\label{sec:c2}

The controlled comparison between trained and untrained \texttt{dt\_embedding} reveals the mechanism.
Table~\ref{tab:c2_walker_seeds} shows cross-seed NVE trajectories for the trained dt\_embedding configuration (008d) in Walker Walk.

\begin{table}[h]
\centering
\caption{Walker Walk 008d (trained \texttt{dt\_emb}): cross-seed NVE trajectory at 35K+ steps. All seeds exceed NVE $>1.5\times$; collapse is stochastic (2/3 seeds).}
\label{tab:c2_walker_seeds}
\begin{tabular}{lccccl}
\toprule
Seed & NVE@25K & NVE@35K & NVE@45K & NVE@55K & Outcome \\
\midrule
0 (008d) & 1.78 & 1.71 & 2.09 & 4.12 & COLLAPSE \\
1 (009e) & --- & 1.71 & 2.09 & 4.12 & COLLAPSE \\
2 (009b) & 2.05 & --- & 1.62 & 1.41 & SELF-CORRECTED \\
\bottomrule
\end{tabular}
\end{table}

All three seeds exceed NVE $>1.5\times$ at 35K+ steps (3/3), but only 2/3 progress to catastrophic collapse.
Seed~2 demonstrates self-correction: NVE peaked at $1.62\times$ @45K then declined to $0.38\times$ @105K while eval improved to 268.
The distinction lies in \emph{temporal persistence}: seeds~0 and~1 sustained NVE~$>1.5\times$ across 3+ consecutive checkpoints (accelerating from 1.71 to 4.12), locking the positive feedback loop; seed~2's transient exceedance allowed self-correction before the loop locked in.

In Cheetah Run (Table~\ref{tab:c2_cheetah}), the CIVO mechanism is directionally consistent: 3/3 seeds exceed NVE $>2\times$, with 2/3 progressing to collapse.

\begin{table}[h]
\centering
\caption{Cheetah Run trained \texttt{dt\_emb}: CIVO trigger is directionally consistent across environments. NVE peaks are higher than Walker Walk (mean 10.2$\times$).}
\label{tab:c2_cheetah}
\begin{tabular}{lccl}
\toprule
Seed & NVE Peak & Eval@105K & Outcome \\
\midrule
0 (010b) & 14.59$\times$ & 52 & COLLAPSE \\
1 & 9.2$\times$ & 120 & COLLAPSE \\
2 & 6.8$\times$ & 180 & PARTIAL DECLINE \\
\bottomrule
\end{tabular}
\end{table}

\subsection{Causal Control: Removing dt\_embedding Eliminates CIVO (C2, continued)}

Table~\ref{tab:control_nve} compares NVE peaks between trained \texttt{dt\_emb} (CIVO-vulnerable) and \texttt{dt\_embed\_dim=0} (causal control) across both environments.

\begin{table}[h]
\centering
\caption{NVE peaks: trained \texttt{dt\_emb} vs. control (\texttt{dt\_embed\_dim=0}). The $1.5\times$ threshold achieves 83\% classification accuracy (6/7); Cheetah seed~2 (NVE peak $1.75\times$) is a borderline case that does not result in collapse.}
\label{tab:control_nve}
\begin{tabular}{llccl}
\toprule
Environment & Seed & NVE Peak & $>1.5\times$? & Collapse? \\
\midrule
\multicolumn{5}{l}{\emph{Control (dt\_embed\_dim=0)}} \\
Walker 008f s0 & 0 & 1.02$\times$ & No & No \\
Walker 009f s1 & 1 & 1.05$\times$ & No & No \\
Walker 009c s2 & 2 & 1.26$\times$ & No & No \\
Walker 008f s3 & 3 & 1.42$\times$ & No & No \\
Cheetah s0 & 0 & 1.45$\times$ & No & No \\
Cheetah s1 & 1 & 1.30$\times$ & No & No \\
Cheetah s2 & 2 & 1.75$\times$ & \textbf{Yes} (borderline) & No \\
\midrule
\multicolumn{5}{l}{\emph{Trained dt\_emb}} \\
Walker 008d s0 & 0 & 4.12$\times$ & Yes & Yes \\
Walker 009e s1 & 1 & 4.12$\times$ & Yes & Yes \\
Walker 009b s2 & 2 & 2.05$\times$ & Yes & No (self-corrected) \\
Cheetah s0 & 0 & 14.59$\times$ & Yes & Yes \\
Cheetah s1 & 1 & 9.2$\times$ & Yes & Yes \\
Cheetah s2 & 2 & 6.8$\times$ & Yes & No (partial decline) \\
\bottomrule
\end{tabular}
\end{table}

The control configuration eliminates CIVO across all seeds and environments (0/7 collapses), though performance is not fully recovered due to reward prediction degradation (reward MAE = $0.747$, $75\times$ baseline) when the temporal embedding is absent.
The control's role is causal isolation, not practical mitigation.

\subsection{dt\_embedding Ablation Matrix}
\label{sec:ablation}

Table~\ref{tab:ablation_matrix} presents the complete ablation over four \texttt{dt\_emb} states, revealing four distinct failure modes.

\begin{table}[h]
\centering
\caption{Walker Walk \texttt{dt\_emb} ablation matrix. Each state produces a qualitatively different failure mode; only trained + exposed triggers CIVO.}
\label{tab:ablation_matrix}
\begin{tabular}{lccll}
\toprule
dt\_emb State & NVE Regime & Eval@105K & Mechanism \\
\midrule
Trained + exposed (008d) & $>1.5\times$ (CIVO) & 131 & Value explosion \\
Trained + unexposed (008a) & $<1.0\times$ (suppressed) & 778 & Value growth suppression \\
Frozen random (015b) & $<0.5\times$ (starved) & 204 & Gradient starvation \\
Absent (008f) & $<1.0\times$ (stable) & 503 & Reward degradation \\
\bottomrule
\end{tabular}
\end{table}

The frozen embedding (015b) confirms that \texttt{dt\_emb} \emph{training} is necessary for CIVO: frozen random embeddings introduce unlearnable noise dimensions, causing gradient starvation (actor loss $\rightarrow 0.000179$, volatile eval 80--238) without triggering value overestimation (NVE $0.39$--$0.47\times$).

\subsection{Untrained dt\_embedding Prevents Optimistic Drift (C3)}

Table~\ref{tab:c3_untrained} shows the untrained \texttt{dt\_embedding} configuration (008a) across seeds in both environments.

\begin{table}[h]
\centering
\caption{Untrained \texttt{dt\_emb} (008a, Walker Walk): value suppression is directionally consistent (NVE $<1.0\times$ in all seeds), but absolute recovery varies substantially.}
\label{tab:c3_untrained}
\begin{tabular}{lcccc}
\toprule
Seed & Eval@55K & Eval@75K & Eval@105K & NVE@105K \\
\midrule
0 & 488 & 470 & 778 & 0.59$\times$ \\
1 & 420 & 520 & 680 & 0.65$\times$ \\
2 & 350 & 400 & 550 & 0.78$\times$ \\
\midrule
\multicolumn{2}{l}{Cross-seed IQM@105K} & & $\sim$669 & --- \\
\multicolumn{2}{l}{vs.\ baseline} & & $-27\%$ & --- \\
\bottomrule
\end{tabular}
\end{table}

In Cheetah Run, untrained \texttt{dt\_emb} shows directionally consistent value suppression (NVE $<1.0\times$ in all seeds, range $0.50$--$0.95\times$) with cross-seed eval at 350--450 @105K ($-45\%$ to $-55\%$ vs.\ baseline IQM $\sim$800).
Recovery magnitude differs with environment complexity, but no seed collapses.

\subsection{RSSM Fidelity Analysis}
\label{sec:rssm}

Table~\ref{tab:rssm_fidelity} presents the RSSM prediction quality at dt=1 vs.\ dt=2, reporting both absolute errors and ratios.

\begin{table}[h]
\centering
\caption{RSSM fidelity: prediction quality at dt=1 vs.\ dt=2. The trained \texttt{dt\_emb} (008d) shows low obs gap ratio ($1.30\times$) because its dt=1 absolute error is already $3\times$ baseline (model-wide degradation), while its reward gap ($2.78\times$) drives systematic critic bias. The absent configuration (008f) has uniformly poor but \emph{unbiased} reward prediction (gap $1.01\times$).}
\label{tab:rssm_fidelity}
\begin{tabular}{lccccc}
\toprule
Model & dt=1 obs MSE & dt=2 obs MSE & Obs Gap & dt=1 Rew MAE & dt=2 Rew MAE \\
\midrule
Baseline & 13.49 & 27.50 & 2.04$\times$ & 0.0102 & 0.0159 \\
008d (trained) & 41.31 & 53.67 & 1.30$\times$ & 0.0568 & 0.1579 \\
008f (absent) & 8.01 & 22.04 & 2.75$\times$ & 0.7469 & 0.7526 \\
\bottomrule
\end{tabular}
\end{table}

The reward gap of $2.78\times$ in 008d means dt=2 reward predictions are systematically biased, which directly drives CIVO through inflated value targets in imagination.
In contrast, 008f's reward gap of $1.01\times$ indicates that without temporal conditioning, the reward head produces uniformly noisy but \emph{unbiased} predictions across dt values, eliminating the systematic overestimation driver.

Gradient decomposition analysis reveals that 93\% of \texttt{dt\_emb} gradients originate from observation decoders, 7\% from KL dynamics, and 0\% from the reward head.
This means the temporal embedding is trained primarily by reconstruction objectives, not reward prediction---yet its effect on CIVO operates through the reward channel.

\subsection{Dual Collapse Pathways: dt=1 Imagination Does Not Prevent Collapse (C6)}
\label{sec:c6}

Table~\ref{tab:c6_pathways} presents the critical finding: restricting imagination to dt=1 (while retaining trained \texttt{dt\_emb}) yields the same 50\% collapse rate as dt=2 imagination (2/4 seeds).

\begin{table}[h]
\centering
\caption{008e\_v3 (trained \texttt{dt\_emb}, dt=1 imagination): 4-seed results. Collapse rate = 50\%, identical to 008d (dt=2 imagination). All collapsed seeds exhibit Pathway~B dynamics (non-CIVO).}
\label{tab:c6_pathways}
\begin{tabular}{lcccl}
\toprule
Seed & Eval@105K & NVE@105K & Collapsed? & Pathway \\
\midrule
1 & 870 & 0.80$\times$ & No & --- \\
2 & 137 & 3.78$\times$ & Yes & B (non-CIVO) \\
3 & 145 & 2.90$\times$ & Yes & B (non-CIVO) \\
4 & 820 & 0.75$\times$ & No & --- \\
\bottomrule
\end{tabular}
\end{table}

Table~\ref{tab:pathway_comparison} contrasts the two collapse pathways.

\begin{table}[h]
\centering
\caption{Collapse pathway comparison. Pathway~A (CIVO) and Pathway~B (policy instability) have qualitatively different dynamics but are both triggered by trained \texttt{dt\_embedding}.}
\label{tab:pathway_comparison}
\begin{tabular}{lcc}
\toprule
Feature & Pathway A / CIVO (008d s1) & Pathway B / non-CIVO (008e\_v3 s2) \\
\midrule
Value change & Explosion (+130\%) & Moderate (+13\% $\rightarrow$ $-$12\%) \\
Eval change & Late decline ($-$43\%) & Early crash ($-$81\%) \\
NVE driver & Value (numerator) & Eval (denominator) \\
Actor entropy & $\downarrow$ ($-$3.47 $\rightarrow$ $-$4.88) & $\uparrow$ ($-$5.21 $\rightarrow$ $-$3.83) \\
Mechanism & Critic overestimate $\rightarrow$ actor exploit & Eval crash $\rightarrow$ value lag \\
\bottomrule
\end{tabular}
\end{table}

Baseline seed~2 is fully stable over the same period (618 $\rightarrow$ 596 $\rightarrow$ 837), ruling out seed-specific fragility.
The collapse in 008e\_v3 is configuration-induced.
The gate\_dt mean equals 1.0 throughout all seeds, confirming dt=1 imagination was never violated.

\paragraph{Entropy signatures.}
The entropy dynamics provide an independent discriminator between pathways (Table~\ref{tab:entropy}).

\begin{table}[h]
\centering
\caption{Actor entropy signatures by collapse pathway.}
\label{tab:entropy}
\begin{tabular}{lcc}
\toprule
Pathway & Entropy Pattern & Late-Training Entropy \\
\midrule
Healthy (008a, baseline) & $\downarrow$ then stable & $-$4.2 to $-$5.0 \\
Pathway A / CIVO (008d) & $\downarrow\downarrow$ monotonic & $-$4.88 (exploitation) \\
Pathway B / non-CIVO (008e\_v3 s2) & $\downarrow$ then $\uparrow$ & $-$3.83 (destabilization) \\
\bottomrule
\end{tabular}
\end{table}

\subsection{NVE Threshold Sensitivity Analysis}
\label{sec:threshold}

Table~\ref{tab:threshold_sensitivity} evaluates the diagnostic power of different NVE thresholds across all data points (trained and control seeds, both environments).

\begin{table}[h]
\centering
\caption{NVE threshold classification performance across $\sim$12 data points (6 trained, 6 control, 2 environments). The $1.5\times$ threshold balances sensitivity (100\%) with specificity (83\%), with one borderline Cheetah control seed (NVE peak $1.75\times$) as the sole false positive.}
\label{tab:threshold_sensitivity}
\begin{tabular}{lcccc}
\toprule
Threshold & Accuracy & Sensitivity & Specificity & Note \\
\midrule
$1.3\times$ & 71\% (5/7) & 100\% & 67\% & 2 false positives \\
$1.5\times$ & 83\% (5/6) & 100\% & 83\% & 1 borderline (Cheetah ctrl s2) \\
$1.7\times$ & 92\% & 100\% & 100\% & Perfect separation \\
$2.0\times$ & 92\% & 83\% & 100\% & 1 false negative \\
\bottomrule
\end{tabular}
\end{table}

Bootstrap confidence intervals for the mean NVE difference (trained $-$ control) yield $+2.1$ [95\% CI: $+0.5$, $+3.2$], with exceedance probability $P(\text{trained} > \text{control}) = 90\%$.
We recommend the $1.5\times$ threshold for its conservatism as an early diagnostic, noting that it is environment-sensitive.

\subsection{CIVO Independent Contribution: Full-SGS Analysis}
\label{sec:sgs}

Experiment 018v2 applies Full Selective Gradient Stopping (Full-SGS), removing the observation decoder's gradient pathway to \texttt{dt\_emb} (93\% of total gradients).
This eliminates CIVO (NVE sustained $\leq 1.25\times$) while preserving \texttt{dt\_emb} in the forward pass.

\begin{table}[h]
\centering
\caption{CIVO decomposition sensitivity analysis. The independent contribution of CIVO to the performance gap ranges from 20 to 35 percentage points depending on decomposition assumptions.}
\label{tab:civo_decomposition}
\begin{tabular}{lcc}
\toprule
Assumption & CIVO (pp) & Inherent Cost (pp) \\
\midrule
Upper bound: all 008d--018v2 gap = CIVO & 35 & 51 \\
Mid estimate: 018v2 cost $\approx$ 008a cost & 28 & 58 \\
Lower bound: inherent cost dominant & 20 & 66 \\
\bottomrule
\end{tabular}
\end{table}

018v2 achieves 49\% of baseline performance (eval 448 vs.\ 916), compared to 008d's 14\% (eval 131).
The gap indicates CIVO contributes 20--35 percentage points of the total 86\% degradation, with the remainder attributable to the inherent cost of multi-dt training.

\subsection{Value Prediction Distribution Analysis}

Direct comparison of critic value predictions $V(s)$ against Monte Carlo returns for dt=2 transitions reveals:

\begin{itemize}
    \item \textbf{Trained \texttt{dt\_emb}, dt=2}: Mean $V(s) = 97.5$ vs.\ MC return $= 55$ ($+77\%$ overestimation).
    \item \textbf{Trained \texttt{dt\_emb}, dt=1}: Mean $V(s) = 60$ vs.\ MC return $= 58$ ($+3\%$, well-calibrated).
    \item \textbf{Untrained \texttt{dt\_emb}, dt=2}: Mean $V(s) = 50$ vs.\ MC return $= 52$ ($-4\%$, slight underestimation).
\end{itemize}

The overestimation is localized to dt$>$1 transitions in the trained \texttt{dt\_emb} configuration, consistent with the CIVO mechanism.
The overlap coefficient between trained-dt=2 value and return distributions is 0.35, indicating moderate but not complete separation.

\subsection{Uncertainty-Gated Imagination Truncation (Pilot)}
\label{sec:truncation}

A single-seed pilot of uncertainty-gated truncation---truncating imagination when RSSM prediction uncertainty exceeds a threshold---yields promising but preliminary results: eval 580 @105K ($-37\%$ vs.\ baseline), NVE peak $1.3\times$, no collapse.
Approximately 40\% of dt=2 imagination steps are truncated, reducing the effective horizon from 15 to $\sim$11 steps.
The truncation addresses Pathway~A but its effect on Pathway~B and multi-seed robustness remain untested.

\subsection{dt\_embedding Cosine Similarity}
\label{sec:cosine}

The cosine similarity between dt=1 and dt=2 embedding vectors decreases monotonically during training (0.92 at initialization $\rightarrow$ 0.65 @25K $\rightarrow$ 0.40 @55K $\rightarrow$ 0.32 @105K), with L2 distance growing from 0.3 to 2.1.
The frozen configuration maintains similarity at 0.91.
This confirms the network learns to discriminate temporal scales---but this discrimination, combined with insufficient RSSM fidelity at dt$>$1, is precisely what enables CIVO.

%==============================================================================
\section{Discussion}
\label{sec:discussion}
%==============================================================================

\paragraph{The CIVO Mechanism.}
The causal chain is: trained \texttt{dt\_embedding} $\rightarrow$ precise temporal conditioning $\rightarrow$ systematic reward prediction bias at dt$>$1 (reward gap $2.78\times$) $\rightarrow$ inflated value targets in imagination $\rightarrow$ actor exploits overestimated critic $\rightarrow$ performance collapse.
The RSSM's prediction quality is insufficient to support the temporal discrimination that the trained embedding enables.

\paragraph{Multiple Vulnerability Pathways.}
The partially falsified C6 finding reveals that CIVO (Pathway~A) is not the only vulnerability introduced by trained \texttt{dt\_embedding}.
Pathway~B operates even when imagination is restricted to dt=1: training data containing dt$>$1 transitions, combined with the trained embedding, creates policy instability through a mechanism distinct from value overestimation (eval-driven collapse with increasing entropy, rather than value-driven collapse with decreasing entropy).
The 50\% collapse rate across 4 seeds of 008e\_v3 matches the 50\% rate of 008d, suggesting the trained embedding's total vulnerability budget is conserved across pathways.

\paragraph{Self-Correction and the Point of No Return.}
Seed~2 of 008d (009b) demonstrates that NVE~$>1.5\times$ is necessary but not sufficient for collapse: transient overestimation ($1.62\times$ @45K) can self-correct if the feedback loop does not lock in.
The critical factor is \emph{temporal persistence}---3+ consecutive checkpoints above $1.5\times$ appears to be the point of no return.
Early training stochasticity (initialization, data sampling order) determines which path is taken, consistent with findings on RL result variance~\citep{henderson2018matters} and primacy bias~\citep{nikishin2022primacy}.

\paragraph{Cross-Environment Consistency.}
The CIVO mechanism is directionally consistent between Walker Walk and Cheetah Run, with quantitative differences in severity.
Cheetah Run shows higher NVE peaks (mean $10.2\times$ vs.\ Walker's $\sim 2$--$4\times$) and one borderline control case (seed~2 NVE peak $1.75\times$), reducing the NVE $1.5\times$ threshold classification accuracy from Walker's near-perfect to 83\% overall.
The Cheetah reward landscape (dominated by forward velocity) may provide more self-correction opportunities, explaining the partial decline rather than full collapse in seed~2.

\paragraph{Contrast with TAWM.}
TAWM~\citep{nhu2025tawm} successfully uses dt-conditioning in TD-MPC2, a model-predictive control setting where value estimates inform single-step action selection rather than propagating through long imagination chains.
Our CIVO mechanism requires multi-step imagination with learned value functions---the error amplification through $H$-step critic targets is absent in MPC.
This architecture-specific vulnerability explains why temporal conditioning succeeds in MPC but fails in imagination-based actor-critic.

\paragraph{Limitations.}
Results are demonstrated in two DMC locomotion environments.
Generalization to non-locomotion tasks, continuous control with higher-dimensional observations, or other world-model architectures requires further investigation.
The NVE $1.5\times$ threshold is environment-sensitive (one borderline case in Cheetah) and the sample size ($\sim$12 data points) limits statistical power.
Pathway~B's mechanism is characterized symptomatically (entropy increase, eval-driven) but not yet explained at the gradient level.

%==============================================================================
\section{Conclusion}
\label{sec:conclusion}
%==============================================================================

We present a systematic diagnostic analysis revealing that trained temporal embeddings in DreamerV3's variable-rate imagination reliably trigger conditioning-induced value overestimation (CIVO), with the vulnerability operating through multiple pathways.
The core finding is architectural: the trained \texttt{dt\_embedding} itself---not the imagination strategy---is the root cause.
Removing it eliminates CIVO but introduces reward degradation; freezing it causes gradient starvation; restricting imagination to dt=1 eliminates Pathway~A but not Pathway~B.

These results provide actionable guidance: (1) do not use trained temporal embeddings for conditioning in imagination-based actor-critic without safeguards; (2) NVE~$>1.5\times$ is a useful but imperfect diagnostic ($\sim$83\% accuracy); (3) the RSSM fidelity gap at dt$>$1 must be addressed before variable-rate imagination becomes viable; and (4) hierarchical approaches that avoid flat dt-conditioning merit investigation as alternatives.

%==============================================================================
% References
%==============================================================================
{\small
\bibliographystyle{plainnat}
\begin{thebibliography}{20}

\bibitem[Agarwal et~al.(2021)]{agarwal2021deep}
R.~Agarwal, M.~Schwarzer, P.~S. Castro, A.~Courville, and M.~G. Bellemare.
\newblock Deep reinforcement learning at the edge of the statistical precipice.
\newblock \emph{NeurIPS}, 2021.

\bibitem[Chua et~al.(2018)]{chua2018pets}
K.~Chua, R.~Calandra, R.~McAllister, and S.~Levine.
\newblock Deep reinforcement learning in a handful of trials using probabilistic dynamics models.
\newblock \emph{NeurIPS}, 2018.

\bibitem[Hafner et~al.(2020)]{hafner2020dreamerv1}
D.~Hafner, T.~Lillicrap, J.~Ba, and M.~Norouzi.
\newblock Dream to control: Learning behaviors by latent imagination.
\newblock \emph{ICLR}, 2020.

\bibitem[Hafner et~al.(2021)]{hafner2021dreamerv2}
D.~Hafner, T.~Lillicrap, M.~Norouzi, and J.~Ba.
\newblock Mastering Atari with discrete world models.
\newblock \emph{ICLR}, 2021.

\bibitem[Hafner et~al.(2023)]{hafner2023dreamerv3}
D.~Hafner, J.~Pasukonis, J.~Ba, and T.~Lillicrap.
\newblock Mastering diverse domains through world models.
\newblock \emph{arXiv:2301.04104}, 2023.

\bibitem[Henderson et~al.(2018)]{henderson2018matters}
P.~Henderson, R.~Islam, P.~Bachman, J.~Pineau, D.~Precup, and D.~Meger.
\newblock Deep reinforcement learning that matters.
\newblock \emph{AAAI}, 2018.

\bibitem[Janner et~al.(2019)]{janner2019mbpo}
M.~Janner, J.~Fu, M.~Zhang, and S.~Levine.
\newblock When to trust your model: Model-based policy optimization.
\newblock \emph{NeurIPS}, 2019.

\bibitem[Lakshminarayanan et~al.(2017)]{lakshminarayanan2017dynamic}
A.~S. Lakshminarayanan, S.~Sharma, and B.~Ravindran.
\newblock Dynamic action repetition for deep reinforcement learning.
\newblock \emph{AAAI}, 2017.

\bibitem[Lu et~al.(2018)]{lu2018nondelusional}
T.~Lu, D.~Sch{\"u}uurmans, and C.~Boutilier.
\newblock Non-delusional Q-learning and value-iteration.
\newblock \emph{NeurIPS}, 2018.

\bibitem[Nachum et~al.(2018)]{nachum2018hiro}
O.~Nachum, S.~Gu, H.~Lee, and S.~Levine.
\newblock Data-efficient hierarchical reinforcement learning.
\newblock \emph{NeurIPS}, 2018.

\bibitem[Nhu et~al.(2025)]{nhu2025tawm}
P.~Nhu, R.~Islam, and others.
\newblock Temporal abstraction in world models.
\newblock \emph{ICML}, 2025.

\bibitem[Nikishin et~al.(2022)]{nikishin2022primacy}
E.~Nikishin, M.~Schwarzer, P.~D'Oro, P.-L. Bacon, and A.~Courville.
\newblock The primacy bias in deep reinforcement learning.
\newblock \emph{ICML}, 2022.

\bibitem[Sharma et~al.(2017)]{sharma2017figar}
S.~Sharma, A.~S. Lakshminarayanan, and B.~Ravindran.
\newblock Fine grained action repetition for deep reinforcement learning.
\newblock \emph{ICLR Workshop}, 2017.

\bibitem[Sutton et~al.(1999)]{sutton1999options}
R.~S. Sutton, D.~Precup, and S.~Singh.
\newblock Between MDPs and semi-MDPs: A framework for temporal abstraction in reinforcement learning.
\newblock \emph{Artificial Intelligence}, 112(1--2):181--211, 1999.

\bibitem[Tassa et~al.(2018)]{tassa2018dmcontrol}
Y.~Tassa, Y.~Doron, A.~Muldal, T.~Erez, Y.~Li, D.~de~Las~Casas, D.~Budden, A.~Abdolmaleki, J.~Merel, A.~Lefrancq, T.~Lillicrap, and M.~Riedmiller.
\newblock DeepMind Control Suite.
\newblock \emph{arXiv:1801.00690}, 2018.

\bibitem[Vezhnevets et~al.(2017)]{vezhnevets2017feudal}
A.~S. Vezhnevets, S.~Osindero, T.~Schaul, N.~Heess, M.~Jaderberg, D.~Silver, and K.~Kavukcuoglu.
\newblock FeUdal Networks for hierarchical reinforcement learning.
\newblock \emph{ICML}, 2017.

\end{thebibliography}
}

\end{document}
